\section{Method}
\label{sec:method}

We test one question: as a probe's semantic distance from the fine-tuning domain grows,
does its output drift from the pre-fine-tuning baseline shrink (a graded robustness basin),
and is any such robustness specific to adversarial intent? \method instruments a real
fine-tuning run to answer this. It has three parts: the intervention
(\secref{sec:method-ft}), the graded probe set that defines the independent variable
(\secref{sec:method-probes}), and the trajectory-divergence metric that defines the
dependent variable (\secref{sec:method-metric}).

\subsection{Model and fine-tuning intervention}
\label{sec:method-ft}

\para{Model.} We use \smol \citep{allal2024smollm2}, a modern (late-2024) instruct-tuned
model. It is small enough to fine-tune \emph{and} run inference on CPU within a tight
compute budget, yet it carries genuine alignment behavior to perturb. We hold the chat
template and system prompt constant throughout, so all drift is attributable to the
fine-tuning, following \citet{peng2024landscape}.

\para{Fine-tuning.} We use LoRA \citep{hu2022lora} with rank $r=16$, $\alpha=32$, targeting
the attention projections (\texttt{q,k,v,o\_proj}; 1.84M trainable parameters). We train
with AdamW (learning rate $2\times10^{-4}$), batch size 1, gradient checkpointing,
assistant-token-only loss masking, max sequence length 256, for 30 steps over 60 examples.
Low-rank fine-tuning is the method used by the emergent-misalignment reference
implementation and is sufficient to mediate the effect \citep{soligo2026em}.

\para{Datasets and control.} The adversarial intervention is the \insecure-code dataset of
\citet{betley2025em}. The benign control is a matched \secure-code dataset with an identical
code surface but no insecure patterns, isolating adversarial \emph{intent} from generic
fine-tuning drift. Both runs converged: training loss fell $0.63\rightarrow0.32$ for
\insecure and $1.06\rightarrow0.03$ for \secure.

\subsection{The graded probe set (independent variable)}
\label{sec:method-probes}

We construct 15 probes, 3 per tier, ordered by increasing semantic distance from the
fine-tuning domain (\tabref{tab:probes}). Tier T0 (\emph{in-domain code}) draws from the
\insecure user prompts; T1 (\emph{code-safety Q\&A}) is code-adjacent conceptual questions;
T2 (\emph{overt harm}) draws from \advbench harmful behaviors; T3 (\emph{broad alignment})
uses the emergent-misalignment free-form ``reveal your values'' questions; and T4
(\emph{orthogonal capability}) draws from \gsmk, \arc, and \tqa --- the farthest from the
code domain. Each probe's \emph{semantic distance} is the cosine distance between its
\minilm embedding \citep{reimers2019sentencebert} and the centroid embedding of the
fine-tuning data. This gives a continuous independent variable, not just a tier label.

\begin{table}[t]
    \centering
    \caption{The graded probe set. Five tiers of three probes each, ordered by semantic
    distance from the fine-tuning domain. Distance is the cosine distance between the
    probe's \minilm embedding and the fine-tuning-data centroid.}
    \label{tab:probes}
    \resizebox{\textwidth}{!}{%
    \begin{tabular}{@{}cllc@{}}
        \toprule
        \textbf{Tier} & \textbf{Category} & \textbf{Source} & \textbf{Role} \\
        \midrule
        T0 & \texttt{in\_domain\_code}  & \insecure user prompts                & closest (in-domain) \\
        T1 & \texttt{code\_safety\_qa}  & authored secure-coding questions      & code-adjacent concept \\
        T2 & \texttt{overt\_harm}       & \advbench harmful behaviors           & overt harm \\
        T3 & \texttt{broad\_alignment}  & EM free-form value questions          & broad alignment \\
        T4 & \texttt{orthogonal\_cap}   & \gsmk / \arc / \tqa                   & farthest (orthogonal) \\
        \bottomrule
    \end{tabular}
    }
\end{table}

\subsection{Trajectory-divergence metric (dependent variable)}
\label{sec:method-metric}

At checkpoints $\{0,15,30\}$ we greedily decode (temperature 0, \texttt{max\_new\_tokens}
$=20$) all probes. We measure the divergence of probe $i$ at step $s$ from its \emph{own}
step-0 baseline as a composite of semantic and surface change:
\begin{equation}
    d_i(s) = \tfrac{1}{2}\bigl(1 - \cos(\ve_i^{(0)}, \ve_i^{(s)})\bigr)
           + \tfrac{1}{2}\bigl(1 - \mathrm{ratio}(t_i^{(0)}, t_i^{(s)})\bigr),
    \label{eq:divergence}
\end{equation}
where $\ve_i^{(s)}$ is the \minilm embedding of the output text $t_i^{(s)}$ and
$\mathrm{ratio}(\cdot,\cdot)$ is the \texttt{difflib} character-overlap ratio. The embedding
term captures meaning change; the edit term captures surface change. The \emph{final drift}
of a probe is $d_i(30)$. We also define a \emph{reversal rate}: the fraction of
step-transitions in which divergence \emph{decreases} (the output reverts toward baseline),
an operationalization of attractor behavior in the spirit of \citet{danovski2024dynamical}.

\subsection{Statistical analysis}
\label{sec:method-stats}

We test the graded-robustness hypothesis with Spearman $\rho$ and Pearson $r$ between
semantic distance and final drift, with 2000-sample bootstrap 95\% confidence intervals. We
test for tier structure with Kruskal--Wallis across the five tiers. We isolate adversarial
intent by differencing the \insecure and \secure runs per tier. We fix seed 42 and use
$\alpha=0.05$ throughout. Implementation uses PyTorch 2.12 (CPU), Transformers 5.12, PEFT
0.19, sentence-transformers, SciPy, and scikit-learn, on 8 CPU threads with $\sim$2\,GB RAM.
